\documentclass[11pt]{article}
\usepackage[margin=2.5cm]{geometry}
\usepackage{amsmath,amssymb}
\usepackage{booktabs}
\usepackage{siunitx}
\usepackage{hyperref}
\usepackage{xcolor}

\sisetup{
  exponent-product = \cdot,
  output-exponent-marker = \ensuremath{\mathrm{e}},
}

\title{Observable Cross-Check of Israel's Approximate\\State-Preparation Circuits}
\author{}
\date{March 2026}

\begin{document}
\maketitle

\section{Introduction}

Israel's circuits are designed to approximate a target quantum state on 72 qubits.
For each circuit, the target state is specified by a set of computational-basis amplitudes stored in a reference file (\texttt{.npz}), containing bitstrings and their associated complex amplitudes $a_{\mathrm{ref}}$.
The circuits themselves are approximate: they are constructed to achieve a target fidelity of roughly $F\approx 0.86$, implying an expected fidelity loss of $\sim 14\%$.

This report describes a local-observable cross-check of the five smallest circuits, \textbf{4q} (4-qubit entanglement clusters, 197 gates), \textbf{6q} (6-qubit clusters, 621 gates), \textbf{8q} (8-qubit clusters, 2336 gates), \textbf{10q} (10-qubit clusters, 7981 gates), and \textbf{12q} (12-qubit clusters, 18376 gates), all acting on 72 qubits.
The goal is to verify that the circuit-output expectation values $\langle Z_i \rangle$ are consistent with those reconstructed from the reference amplitudes, within the bounds implied by the known fidelity.

\section{Methods}

\subsection{Simulator 1: Pauli Propagation (exact)}

We use \textsc{PauliPropagation.jl} to compute exact circuit-output expectation values in the Heisenberg picture.
For each qubit~$i$, the observable $Z_i$ is propagated backward through the circuit:
\begin{equation}
  \langle Z_i \rangle_{\mathrm{circuit}} = \langle 0^{\otimes n} | \, U^\dagger Z_i \, U \, | 0^{\otimes n} \rangle
  = \texttt{overlapwithzero}\!\bigl(\texttt{propagate}(\text{circuit},\, Z_i)\bigr).
\end{equation}
The initial state is $|0\rangle^{\otimes 72}$.
All truncation is disabled: \texttt{min\_abs\_coeff}${}=0$, \texttt{max\_weight}${}=\infty$, ensuring the result is exact to machine precision.

Each Qiskit $U(\theta,\phi,\lambda)$ gate is decomposed as $R_z(\phi)\,R_y(\theta)\,R_z(\lambda)$, and each \texttt{cx} gate maps to a Clifford CNOT.

\subsection{Simulator 2: TEBD (MPS-based)}

As an independent cross-check, we simulate the same circuits using TEBD (time-evolving block decimation) within our MPS framework.
The initial state is $|0\rangle^{\otimes 72}$ in MPS form.
For the 4q circuit, TEBD runs without truncation (maximum bond dimension reached is 83).
For the 6q circuit, the bond dimension saturates at the cap of $\chi_{\max}=512$ from layer~14 onward (out of 44 layers), introducing significant truncation error.
For the larger 8q, 10q, and 12q circuits, no TEBD results are included here: this report extends those cases only with exact Pauli~Propagation, since MPS simulation becomes substantially more expensive and was deferred.

\subsection{Reference observables from target amplitudes}

The target-state expectation values are reconstructed from the reference data as:
\begin{equation}
  \langle Z_i \rangle_{\mathrm{target}} = \sum_{\mathbf{b}} |a_{\mathrm{ref}}(\mathbf{b})|^2 \, (-1)^{b_i},
\end{equation}
where the sum runs over all stored bitstrings $\mathbf{b}$ and $b_i$ is the bit at qubit~$i$.
The bitstrings use Qiskit big-endian ordering, where the leftmost character corresponds to the highest-indexed qubit.

\subsection{Validation of the gate mapping}

The circuit conversion was validated on small hand-checkable circuits by comparing Pauli~Propagation against Qiskit's statevector simulator.
For a single $U(\theta,\phi,\lambda)$ gate with random parameters, all three Pauli expectations $\langle X\rangle$, $\langle Y\rangle$, $\langle Z\rangle$ agree with Qiskit to machine precision ($< 10^{-15}$).
For two-qubit circuits involving $U$ gates followed by CNOT (both control/target orderings), $\langle Z_i\rangle$ values agree with Qiskit to $< 10^{-15}$.

\section{Results}

\subsection{Pauli Propagation vs.\ TEBD on the 4q circuit}

For the 4q circuit, which does not require bond-dimension truncation, Pauli~Propagation and TEBD agree on all 72 site-resolved $\langle Z_i \rangle$ values to machine precision:
\begin{center}
\begin{tabular}{lS[table-format=1.2e-1]S[table-format=1.2e-1]}
\toprule
 & {MAE} & {Max error} \\
\midrule
All qubits (72)     & 2.87e-12 & 1.32e-11 \\
Active qubits (36)  & 5.74e-12 & 1.32e-11 \\
Inactive qubits (36)& 3.12e-16 & 7.77e-16 \\
\bottomrule
\end{tabular}
\end{center}
This confirms that both simulators produce identical circuit-output observables.

\subsection{Pauli Propagation vs.\ TEBD on the 6q circuit}

For the 6q circuit, TEBD saturates the bond dimension at $\chi=512$ and produces unreliable results (the TEBD--reference amplitude overlap gives a fidelity of $F\approx 0.0003$, far from the expected $F\approx 0.86$).
Pauli~Propagation, being exact, is the reliable reference here.
On inactive qubits (which involve only single-qubit gates and are unaffected by bond truncation), the two methods still agree to $\sim 10^{-14}$.

\subsection{Circuit output vs.\ target state}

Table~\ref{tab:main} summarises the comparison of exact circuit-output $\langle Z_i \rangle$ (from Pauli~Propagation) against the target-state $\langle Z_i \rangle$ (from the reference amplitudes).

\begin{table}[h]
\centering
\caption{Comparison of circuit-output and target-state $\langle Z_i \rangle$ for the 4q, 6q, 8q, 10q, and 12q circuits.
Active qubits are those involved in at least one CNOT gate; inactive qubits undergo only single-qubit rotations.}
\label{tab:main}
\begin{tabular}{llS[table-format=1.2e-1]S[table-format=1.2e-1]S[table-format=1.2e-1]}
\toprule
Circuit & Qubit group & {MAE} & {Max $|\Delta Z_i|$} & {RMSE} \\
\midrule
4q & All (72)       & 1.55e-2 & 4.76e-2 & 2.27e-2 \\
   & Active (36)    & 2.65e-2 & 4.76e-2 &         \\
   & Inactive (36)  & 4.42e-3 & 2.94e-2 &         \\
\midrule
6q & All (72)       & 1.46e-2 & 5.04e-2 & 2.12e-2 \\
   & Active (40)    & 2.42e-2 & 5.04e-2 &         \\
   & Inactive (32)  & 2.54e-3 & 6.90e-3 &         \\
\midrule
8q & All (72)       & 1.30e-2 & 3.99e-2 & 1.86e-2 \\
   & Active (40)    & 2.21e-2 & 3.99e-2 &         \\
   & Inactive (32)  & 1.61e-3 & 8.52e-3 &         \\
\midrule
10q & All (72)      & 1.20e-2 & 3.80e-2 & 1.68e-2 \\
    & Active (50)   & 1.72e-2 & 3.80e-2 &         \\
    & Inactive (22) & 1.01e-4 & 1.32e-3 &         \\
\midrule
12q & All (72)      & 1.15e-2 & 4.00e-2 & 1.61e-2 \\
    & Active (48)   & 1.72e-2 & 4.00e-2 &         \\
    & Inactive (24) & 2.02e-4 & 3.64e-3 &         \\
\bottomrule
\end{tabular}
\end{table}

For the larger 8q, 10q, and 12q circuits, exact Pauli~Propagation continues to show the same qualitative pattern as for 4q and 6q: the dominant discrepancies are confined to active qubits, while inactive qubits remain extremely close to the target-state values.
The overall errors decrease slightly with circuit size in this set, from MAE $\approx 1.55\times 10^{-2}$ for 4q to $1.15\times 10^{-2}$ for 12q.

\subsection{Fidelity bound}

For two states $|\psi\rangle$ and $|\phi\rangle$ with fidelity $F = |\langle \psi | \phi \rangle|^2$, the difference in expectation value of any observable with $\|O\| \leq 1$ satisfies
\begin{equation}
  |\langle O \rangle_\psi - \langle O \rangle_\phi| \leq 2\bigl(1 - \sqrt{F}\bigr).
\end{equation}
For the 4q circuit, the fidelity between the TEBD circuit output and the reference state was computed directly from the stored amplitudes:
\begin{equation}
  F_{4\mathrm{q}} = \bigl|\textstyle\sum_{\mathbf{b}} a_{\mathrm{ref}}^*(\mathbf{b})\, a_{\mathrm{circuit}}(\mathbf{b})\bigr|^2 = 0.8590,
\end{equation}
matching the expected value.
The resulting fidelity bound is $2(1-\sqrt{0.859}) \approx 0.146$, and the observed maximum $|\Delta Z_i| = 0.048$ is well within this bound.

For the 6q circuit, the direct fidelity computation from TEBD is unreliable due to bond-dimension truncation.
However, the observed $\langle Z_i \rangle$ errors from Pauli~Propagation (max $|\Delta Z_i| = 0.050$) are consistent with the reported fidelity $F_{6\mathrm{q}} \approx 0.865$ and bound $2(1-\sqrt{0.865}) \approx 0.141$.

For the 8q, 10q, and 12q circuits, we use the expected fidelities reported with the circuit data:
\begin{align}
  F_{8\mathrm{q}} &\approx 0.8794, & 2(1-\sqrt{F_{8\mathrm{q}}}) &\approx 0.124, \\
  F_{10\mathrm{q}} &\approx 0.8866, & 2(1-\sqrt{F_{10\mathrm{q}}}) &\approx 0.117, \\
  F_{12\mathrm{q}} &\approx 0.8890, & 2(1-\sqrt{F_{12\mathrm{q}}}) &\approx 0.114.
\end{align}
The observed maxima, $|\Delta Z_i|_{\max} = 0.040$ for 8q, $0.038$ for 10q, and $0.040$ for 12q, remain comfortably within these bounds.

\subsection{Bitstring convention}

Two bitstring ordering conventions were tested for reconstructing $\langle Z_i\rangle_{\mathrm{target}}$:
\begin{center}
\begin{tabular}{llS[table-format=1.2e-1]S[table-format=1.1]}
\toprule
Convention & Description & {MAE (4q)} & {MAE (native)} \\
\midrule
Qiskit big-endian & leftmost bit = highest qubit & 1.55e-2 & \\
Native (little-endian) & leftmost bit = lowest qubit & & 1.3 \\
\bottomrule
\end{tabular}
\end{center}
The Qiskit big-endian convention produces errors consistent with the fidelity bound; the native convention yields errors close to~2 (effectively random), confirming that the reference data uses Qiskit ordering.
The same convention was selected automatically for the 8q, 10q, and 12q runs as well, with native-ordering MAE again of order unity ($\sim 1.33$), while the Qiskit-ordering MAE stays at the $\sim 10^{-2}$ level.

\subsection{Single-qubit X/Y check on the 4q circuit}

As an additional phase-sensitive validation, we compared exact Pauli~Propagation against the sparse reference amplitudes for all 72 single-qubit observables $X_i$, $Y_i$, and $Z_i$ on the 4q circuit.
With the corrected gate mapping, the same Qiskit-consistent bit ordering (implemented as the reversed convention in the local sparse-state parser) is selected by a clear margin.

For the single-qubit off-diagonal observables, all reference expectations vanish within numerical precision:
\begin{center}
\begin{tabular}{lS[table-format=1.2e-1]S[table-format=1.2e-1]}
\toprule
Observable family & {MAE} & {Max error} \\
\midrule
$X_i$ (72 qubits) & 2.63e-16 & 1.89e-15 \\
$Y_i$ (72 qubits) & 2.11e-16 & 1.38e-15 \\
\bottomrule
\end{tabular}
\end{center}
Pauli~Propagation gives the same result: all 72 single-qubit $X_i$ and $Y_i$ expectations are zero up to machine precision.
This makes the check reliable in the limited sense that there is no detectable phase-convention mismatch between the sparse reference amplitudes and Pauli~Propagation for these observables.
However, it is not a strong off-diagonal stress test, because the signal is trivial: all single-qubit $X_i$ and $Y_i$ expectations are essentially zero.
To probe genuinely nontrivial off-diagonal structure, two-qubit correlators such as $X_iX_j$ or $Y_iY_j$ would be more informative.

\section{Discussion}

\begin{enumerate}
  \item \textbf{Gate mapping is validated.}
    Pauli Propagation matches Qiskit's statevector simulator to machine precision on small test circuits and agrees with TEBD to $\sim 10^{-12}$ on the full 72-qubit 4q circuit.

  \item \textbf{The reference data uses Qiskit big-endian bitstring ordering.}
    This was confirmed by testing both plausible conventions: only the big-endian convention gives physically consistent $\langle Z_i\rangle$ values.
    The dedicated 4q single-qubit $X_i/Y_i$ check is also consistent with this convention.

  \item \textbf{All $\langle Z_i \rangle$ discrepancies are within the fidelity bound.}
    Across 4q, 6q, 8q, 10q, and 12q, the maximum observed error ranges from $\sim 0.038$ to $\sim 0.050$, well within the theoretical bounds of $\sim 0.11$--$0.15$ implied by the reported fidelities.
    The discrepancies are concentrated on active qubits (those involved in CNOT gates), consistent with the fact that the circuit approximation error is carried by the entangling part of the circuit.

  \item \textbf{Single-qubit $X_i/Y_i$ checks on 4q are consistent but not discriminating.}
    Both the reference amplitudes and Pauli~Propagation give $\langle X_i\rangle \approx \langle Y_i\rangle \approx 0$ for all 72 qubits, at the level of machine precision.
    This supports the correctness of the off-diagonal observable evaluation, but because the expected signal vanishes, it does not by itself provide a stringent test of the phase structure.

  \item \textbf{TEBD at $\chi=512$ is insufficient for the 6q circuit.}
    The 6q circuit requires bond dimensions exceeding 512 for accurate simulation.
    Pauli~Propagation, which is exact for observable-level computations, is the appropriate tool for this cross-check and was therefore used as the sole method extended to 8q, 10q, and 12q here.

  \item \textbf{The circuits behave as expected.}
    There is no evidence of an anomaly in the circuit output.
    The local observable values produced by the circuits are consistent with approximate preparation of the target state at the reported fidelity.
\end{enumerate}

\end{document}
